\section{Results}
\label{sec:results}

This section presents out-of-sample forecasting results from two of the three analyses described in Section~\ref{sec:experimental_design}: the Ridge--HAR and Ridge--PCA subgroup analyses. All metrics are computed on strictly out-of-sample predictions generated by the walk-forward procedure of Section~\ref{sec:backtest}. MSE and MAE are evaluated on the preprocessed (adjusted) scale, while QLIKE is computed on the raw scale after applying Duan's smearing estimator (Section~\ref{sec:duan}).

\subsection{Ridge--HAR Subgroup Analysis}
\label{sec:results_ridge_har}

Table~\ref{tab:ridge_har_subgroup} reports the forecasting performance of ridge regression with HAR features across all exogenous variable subgroups. By holding the model fixed and varying the information set, this analysis isolates the marginal predictive content of each subgroup beyond the target's own lagged history.

\begin{table}[H]
\centering
\caption{Out-of-sample performance of Ridge--HAR across exogenous variable subgroups. $\Delta$QLIKE denotes the change relative to the baseline (negative $=$ improvement). OOS~$R^2$ is relative to the Ridge--HAR baseline. $N$ denotes the number of out-of-sample observations.}
\label{tab:ridge_har_subgroup}
\begin{tabular}{lrccccr}
\toprule
Subgroup & $N$ & MSE & MAE & QLIKE & $\Delta$QLIKE & OOS $R^2$ (\%) \\
\midrule
Baseline         & 222{,}058 & 0.0629 & 0.1870 & 0.1380 & ---      & 0.0 \\
Moments          & 217{,}214 & 0.0617 & 0.1851 & 0.1322 & $-$0.0059 & 2.0 \\
Liquidity        & 203{,}433 & 0.0620 & 0.1856 & 0.1317 & $-$0.0063 & 1.5 \\
Market (EW)      & 200{,}687 & 0.0622 & 0.1861 & 0.1338 & $-$0.0043 & 1.2 \\
Market (VW)      & 200{,}687 & 0.0621 & 0.1859 & 0.1335 & $-$0.0046 & 1.4 \\
Sentiment        & 138{,}279 & 0.0652 & 0.1881 & 0.1397 & $+$0.0016 & $-$3.6 \\
Implied Vol      &  14{,}468 & 0.0595 & 0.1782 & 0.1259 & $-$0.0121 & 5.5 \\
Vol Demand       & 136{,}849 & 0.0649 & 0.1877 & 0.1405 & $+$0.0025 & $-$3.1 \\
\textbf{All Features} & \textbf{12{,}890} & \textbf{0.0587} & \textbf{0.1758} & \textbf{0.1231} & $\boldsymbol{-}$\textbf{0.0149} & \textbf{6.7} \\
\bottomrule
\end{tabular}
\end{table}

The Ridge--HAR baseline, which uses only the target's own HAR features, already achieves strong forecasting accuracy (MSE $= 0.0629$, QLIKE $= 0.1380$). The geometric lag structure of Section~\ref{sec:har_features} effectively encodes multi-scale volatility persistence, leaving limited room for marginal improvement from exogenous variables.

Results fall into three tiers. First, implied volatility and the combined feature set deliver meaningful improvements: implied volatility reduces QLIKE by 8.8\% (OOS~$R^2 = 5.5\%$), while the all-features configuration achieves the largest gains across all metrics (OOS~$R^2 = 6.7\%$, $\Delta$QLIKE $= -0.0149$). Second, moments, liquidity, and market factors provide modest improvements in the range of OOS~$R^2 = 1$--$2\%$. Third, sentiment and volatility demand variables \emph{degrade} performance relative to the baseline (OOS~$R^2 = -3.6\%$ and $-3.1\%$, respectively).

These results should be interpreted with attention to sample sizes. The implied volatility and all-features subgroups have substantially fewer out-of-sample observations ($N \approx 13$--$14\text{K}$ vs.\ $222\text{K}$ for the baseline) due to shorter data availability windows for options-derived variables. Sentiment and volatility demand also operate on reduced samples ($N \approx 137$--$138\text{K}$), covering different time periods that may encompass distinct volatility regimes. The strong performance of options-implied information, while consistent with the forward-looking nature of these variables, requires confirmation on longer samples.

\subsection{HAR Model Comparison}
\label{sec:results_model_comparison}

% TODO: Results pending. Compare naive, ridge, XGBoost, LightGBM, and random forest
% on baseline HAR features (no exogenous variables).

\subsection{Ridge--PCA Subgroup Analysis}
\label{sec:results_ridge_pca}

Table~\ref{tab:ridge_pca_subgroup} reports results for ridge regression with PCA features. The PCA approach replaces the hand-crafted HAR lag structure with data-driven dimensionality reduction (Section~\ref{sec:pca_features}), and OOS~$R^2$ is computed relative to the naive baseline.

\begin{table}[H]
\centering
\caption{Out-of-sample performance of Ridge--PCA across exogenous variable subgroups. $\Delta$QLIKE denotes the change relative to the naive baseline (negative $=$ improvement). OOS~$R^2$ is relative to the naive baseline. $N$ denotes the number of out-of-sample observations.}
\label{tab:ridge_pca_subgroup}
\begin{tabular}{lrccccr}
\toprule
Subgroup & $N$ & MSE & MAE & QLIKE & $\Delta$QLIKE & OOS $R^2$ (\%) \\
\midrule
Naive Baseline   & 221{,}915 & 0.2248 & 0.3431 & 0.3518 & ---      & 0.0 \\
PCA Baseline     & 218{,}934 & 0.1337 & 0.2644 & 0.2621 & $-$0.0897 & 40.5 \\
\textbf{Moments} & \textbf{214{,}090} & \textbf{0.0722} & \textbf{0.2001} & \textbf{0.1497} & $\boldsymbol{-}$\textbf{0.2021} & \textbf{67.9} \\
Liquidity        & 200{,}309 & 0.1300 & 0.2588 & 0.2509 & $-$0.1008 & 42.2 \\
Market (EW)      & 196{,}753 & 0.1132 & 0.2456 & 0.2212 & $-$0.1306 & 49.7 \\
Market (VW)      & 196{,}753 & 0.1082 & 0.2408 & 0.2130 & $-$0.1388 & 51.9 \\
Sentiment        & 135{,}155 & 0.0839 & 0.2149 & 0.1692 & $-$0.1826 & 62.7 \\
Implied Vol      &  11{,}344 & 0.1652 & 0.3009 & 0.2900 & $-$0.0618 & 26.5 \\
Vol Demand       & 131{,}172 & 0.0827 & 0.2128 & 0.1687 & $-$0.1831 & 63.2 \\
All Features     &   9{,}766 & 0.0988 & 0.2381 & 0.1957 & $-$0.1561 & 56.1 \\
\bottomrule
\end{tabular}
\end{table}

The Ridge--PCA baseline achieves an OOS~$R^2$ of 40.5\% over the naive forecast, confirming that PCA successfully extracts predictive lag structure from the raw data. However, its absolute MSE ($0.1337$) is more than double that of the Ridge--HAR baseline ($0.0629$), indicating that five PCA components estimated on a 500-observation rolling window cannot match the hand-crafted HAR rolling-mean structure for encoding volatility persistence.

The most striking result is the magnitude of marginal gains from exogenous variables under PCA, which far exceed those observed under HAR. Moments achieve OOS~$R^2 = 67.9\%$, followed by volatility demand ($63.2\%$) and sentiment ($62.7\%$). This contrast is especially sharp for sentiment and volatility demand, which \emph{degrade} Ridge--HAR performance but provide large improvements under PCA. The explanation is that the HAR lag structure already encodes the dominant predictive signal---multi-scale volatility persistence---leaving exogenous variables to contribute only incremental (and sometimes noisy) information. The weaker PCA baseline, by contrast, leaves substantial room for exogenous variables to compensate for what the data-driven projection misses. In particular, moment variables (bipower variation, semi-variance, higher-order returns) carry strong persistence information that HAR features capture by construction but PCA features do not.

As with the HAR analysis, sample size discrepancies warrant caution. Implied volatility ($N = 11{,}344$) and all features ($N = 9{,}766$) have an order of magnitude fewer observations than the baseline, and their relatively weaker OOS~$R^2$ may partly reflect overfitting in the high-dimensional, small-sample regime despite ridge regularization.

\subsection{HAR vs.\ PCA: Comparative Assessment}
\label{sec:results_comparison}

Comparing across the two feature engineering strategies reveals a clear trade-off. Ridge--HAR dominates in absolute forecasting accuracy: even the HAR baseline (MSE $= 0.0629$) outperforms every PCA configuration except moments (MSE $= 0.0722$), and the best HAR configuration---all features---achieves the lowest MSE overall ($0.0587$). The hand-crafted geometric lag structure of Section~\ref{sec:har_features}, informed by the known persistence properties of realized volatility, provides a structural advantage that data-driven PCA compression does not overcome.

Nevertheless, the PCA analysis serves as a valuable diagnostic. The large marginal gains from exogenous variables under PCA reveal which information sources carry the most volatility-relevant content: higher-order return moments, volatility demand proxies, and social media sentiment all contribute substantial predictive power when the baseline feature set is weaker. That these same variables add little (or even hurt) under HAR confirms that the HAR lag structure already captures much of the information they provide. The practical implication is that Ridge--HAR augmented with options-implied variables represents the most promising configuration, achieving the best absolute performance while incorporating genuinely complementary forward-looking information---though this conclusion rests on a limited sample ($N = 12{,}890$) and requires validation on longer options data histories.
